\documentclass[fontset=windows]{article}
\usepackage[margin=1in]{geometry}%设置边距，符合Word设定
\usepackage{fontspec}
\setmainfont{Calibri}
\usepackage{ctex}
\usepackage{setspace}
\usepackage{lipsum}
\usepackage{appendix}
\usepackage{booktabs}	% 加入表格
\usepackage{amsmath}    % 数学
\numberwithin{equation}{section}
\usepackage{amssymb}    % 数学符号
\usepackage{amsthm}     % 定理环境
\usepackage{bm}         % 数学粗体
\usepackage{tikz}       % 绘图
\usepackage{colortbl}	% 高亮表格
\usepackage[dvipsnames]{xcolor} 	% 字体颜色
\usepackage{graphicx}	% 加入图片
\usetikzlibrary{patterns}
\usepackage[shortlabels,inline]{enumitem}
\usepackage{caption}
\usepackage{subcaption}
\graphicspath{{Figures/}} 
\captionsetup[figure]{
    name=Fig.,
    labelsep=period,
    font=small,
    labelfont=bf
}
\numberwithin{figure}{section}
\usepackage{hyperref} % 对目录生成链接，注：该宏包可能与其他宏包冲突，故放在所有引用的宏包之后
\usepackage[nameinlink]{cleveref} % 更智能的交叉引用
\hypersetup{colorlinks = true,  % 将链接文字带颜色
	bookmarksopen = true, % 展开书签
	bookmarksnumbered = true, % 书签带章节编号
	pdftitle = Individual Decision Making, % 标题
	pdfauthor = Zhuangfu Zhao} % 作者
\bibliographystyle{plain} % 参考文献引用格式
\newcommand{\upcite}[1]{\textsuperscript{\cite{#1}}}
\renewcommand{\contentsname}{\centerline{Contents}} %经过设置word格式后，将目录标题居中

%---------- 定理设置  --------------
\theoremstyle{plain}
%\addtobeamertemplate{theorem begin}{\normalfont\rmfamily}{}
\newtheorem{theorem}{Theorem}
\numberwithin{theorem}{section}
\newtheorem{lemma}{Lemma}
\numberwithin{lemma}{section}
\newtheorem{corollary}{Corollary}
\numberwithin{corollary}{section}
%\theoremstyle{definition}
\newtheorem{definition}{Definition}
\numberwithin{definition}{section}
\newtheorem{proposition}{Proposition}
\numberwithin{proposition}{section}
\newtheorem{example}{Example}
\numberwithin{example}{section}
\newtheorem{assumption}{Assumption}
\numberwithin{assumption}{section}

%---------- 证明和解环境 ----------
\renewenvironment{proof}[1][Proof]{\par\textbf{#1:}%
\enspace\ignorespaces}{\par}
\newenvironment{solution}{\par\textbf{Solution:}%
\enspace\ignorespaces}{\par}

%---------- 数学符号 ----------
\newcommand{\ud}{\mathrm{d}}

%---------- 证明符号 ----------
\newcommand{\QED}{\hfill\ensuremath{\square}}
\newcommand{\logiceq}{\Leftrightarrow}
\newcommand{\suff}{\Rightarrow}
\newcommand{\nece}{\Leftarrow}

%---------- 环境美化 ----------
% \usepackage[most]{tcolorbox}
% \newcommand{\newtcbenvironment}[2]{
%     \tcolorboxenvironment{#1}{#2, enhanced, breakable, sharp corners, boxrule=1pt}
%     \tcolorboxenvironment{#1*}{#2, enhanced, breakable, rounded corners, boxrule=1pt}
% }
% \newtcbenvironment{theorem}{colframe=RoyalPurple, colback=RoyalPurple!8}
% \newtcbenvironment{proposition}{colframe=RoyalPurple, colback=RoyalPurple!8}
% \newtcbenvironment{corollary}{colframe=NavyBlue, colback=SkyBlue!8}
% \newtcbenvironment{lemma}{colframe=NavyBlue, colback=SkyBlue!8}
% \newtcbenvironment{claim}{colframe=NavyBlue, colback=SkyBlue!8}

% \newtcbenvironment{definition}{colframe=ForestGreen, colback=ForestGreen!5}
% \newtcbenvironment{example}{colframe=RawSienna, colback=RawSienna!5}
% \newtcbenvironment{problem}{colframe=WildStrawberry!30, colback=WildStrawberry!5}

%---------- 调节公式的间距 ----------
%\AtBeginDocument{
%	\setlength{\abovedisplayskip}{4pt plus 1pt minus 1pt}
%	\setlength{\belowdisplayskip}{4pt plus 1pt minus 1pt}
%	\setlength{\abovedisplayshortskip}{2pt}
%	\setlength{\belowdisplayshortskip}{2pt}
%	\setlength{\arraycolsep}{2pt}
%}


\title{\heiti\zihao{2} Quantitative Spatial Model: Theory and Application}
\author{\songti Tianyi Liang and Zhuangfu Zhao}
\date{\today}


\begin{document}
\maketitle
\thispagestyle{empty}

\setcounter{tocdepth}{2}
\tableofcontents


\section{Introduction}

\noindent Classical trade theories (Ricardo, Heckscher-Ohlin), while extremely useful in highlighting the economic forces behind trade, 
are very difficult to generalize to a set-up with many trading partners and bilateral trade costs. 
Thus the classical theories do not provide much guidance in doing empirical work (lack of quantifiability). Because of this difficulty, 
those doing empirical work in trade began using a statistical model known as the gravity equation. 

It is no exaggeration to say that the gravity equation is one of the most successful empirical relationships in all of economics. 
Despite the gravity equation's empirical success, its lack of theoretical foundations and interpretability long limited its use for policy analysis or counterfactual predictions. 
This gap between empirical relevance and theoretical rigor motivated the development of new models that could both explain the gravity relationship and provide a structural basis. 

The Eaton-Kortum (henceforth EK) model, rooted in the Ricardian tradition, addresses this need by introducing a probabilistic representation of productivity differences across countries and goods. 
It not only provides a micro-founded derivation of the gravity equation but also incorporates bilateral trade costs and multiple countries naturally, 
making it fully consistent with real-world trade patterns. We focus on the evolution of international trade models from the Ricardo model to 
the EK model, followed by a replication to Tombe \& Zhu (2019).

\section{The Ricardo Model}




\end{document}
